(1) 对容器内 Maxwell 分布，法向正速度分量 $v_z>0$ 的粒子穿过面积 $A$ 的通量为
\[
-\dot N=A\frac NV\int_{v_z>0}v_z f(\boldsymbol v)\,d^3v
=A\frac NV\sqrt{\frac{k_BT}{2\pi m}}.
\]
由于穿孔权重多乘一个 $v_z$，逸出者比均匀抽取的一颗气体分子更可能具有较大法向速率。对通量加权后的动能积分：
\[
\langle\varepsilon\rangle_{\rm out}
=\frac{\int_{v_z>0}\frac12mv^2\,v_z f\,d^3v}
{\int_{v_z>0}v_z f\,d^3v}
=2k_BT,
\]
而容器内一颗粒子的平均动能是 $3k_BT/2$。

(2) 容器内总内能 $U=3Nk_BT/2$。每减少 $-dN$ 颗粒子（$dN<0$），离开容器的能量为 $(-dN)\,2k_BT$，因此
\[
dU=2k_BT\,dN.
\]
将 $U$ 微分得
\[
\frac32k_B(N\,dT+T\,dN)=2k_BT\,dN
\quad\Longrightarrow\quad \frac{dT}{T}=\frac13\frac{dN}{N}.
\]
积分得到
\[
T=T_0(N/N_0)^{1/3}.
\]
气体逸出时 $N$ 下降、$T$ 也下降，这种冷却不需要对外做机械活塞功。

(3) 定义初始逸出时间尺度
\[
\tau_0=\frac{V}{A}\sqrt{\frac{2\pi m}{k_BT_0}}.
\]
将 $T(N)$ 代回逸出率方程，令 $x=N/N_0$，则
\[
\dot x=-\frac1{\tau_0}x^{7/6}.
\]
求解给 $x^{-1/6}=1+t/(6\tau_0)$，于是
\[
N(t)=N_0\left(1+\frac{t}{6\tau_0}\right)^{-6},\qquad
T(t)=T_0\left(1+\frac{t}{6\tau_0}\right)^{-2}.
\]
因 $p=Nk_BT/V$，
\[
p(t)=p_0\left(1+\frac{t}{6\tau_0}\right)^{-8}.
\]
代入 $t=0$ 和原速率方程，可独立核查初值及幂指数。

(4) 容器内随机粒子平均动能为 $3k_BT/2$，但开孔造成速度选择：穿孔概率正比正法向速率，从而每个逸出粒子携带 $2k_BT$，比容器平均更高的能量，因此余气冷却。上述代数长尾只在 $N\gg1$、孔面积远小于壁面、流动足够缓慢使容器内部持续热平衡时成立；当粒子太稀少导致内部碰撞不足时，Maxwell 假设失效，不能将结果延伸到任意晚时。
